\subsection{SUBSPACES OF COMPLETE SPACES}

PROPOSITION 8. Every closed subspace of a complete space is complete. Every complete subspace of a Hausdorff uniform space (complete or not) is closed.

Let X be a complete space and let A be a closed subspace of X. If F is a Cauchy filter on A, then it is a Cauchy filter base on X (no. 1, Proposition 3) and therefore converges to a point \(x \in X\) ; but since A is closed we have \(x \in A\) , and therefore F converges in the subspace A.

Now let A be a non-closed subset of a Hausdorff uniform space X, and let \(b \in \overline{\mathbf{A}} - \mathbf{A}\) . The trace \(\mathfrak{B}_{\mathbf{A}}\) on A of the neighbourhood filter B of b in X is a Cauchy filter on A; but it cannot converge to a point \(c \in \mathbf{A}\) , otherwise c would be a limit point of B (no. 2, Proposition 5, Corollary 3) which is absurd since \(b \neq c\) and X is Hausdorff.

PROPOSITION 9. Let X be a uniform space and let A be a dense subset of X such that every Cauchy filter base on A converges in X. Then X is complete.

It is enough to show that every minimal Cauchy filter \(\mathfrak{F}\) on X is convergent. Since A is dense and since every set of \(\mathfrak{F}\) has a non-empty interior (no. 2, Corollary 4 to Proposition 5), the trace \(\mathfrak{F}_{\mathrm{A}}\) of \(\mathfrak{F}\) on A is a Cauchy filter on A, hence converges to a point \(x_{0} \in \mathrm{X}\) . Since \(\mathfrak{F}\) is coarser than the filter on X generated by \(\mathfrak{F}_{\mathrm{A}}\) , it follows that \(\mathfrak{F}\) converges to \(x_{0}\) (no. 2, Corollary 3 to Proposition 5).

